\documentclass[11pt]{article}
\usepackage[utf8]{inputenc}
\usepackage{amsmath, amssymb, amsthm}
\usepackage{geometry}
\geometry{letterpaper, margin=1in}
\usepackage{hyperref}
\usepackage{enumitem}

\title{Homework 4.10: Forced Mechanical Vibrations \\ Problem Sets and Solutions}
\author{Curated Practice Sets}
\date{}

\begin{document}

\maketitle
\tableofcontents
\newpage

% ---------------------------------------------------------
% Concept Checklist & Question Mapping
% ---------------------------------------------------------
\section{Concept Checklist \& Question Mapping}
To ensure comprehensive coverage of all topics, tricks, and problem archetypes found in the source material, the following concepts have been identified and mapped to the generated problems:

\begin{itemize}
    \item \textbf{C1: Frequency Response Curve \& Damped Resonance} (Checking $b^2 < 2mk$, calculating $\gamma_r$ and maximum gain $M(\gamma_r)$). \textit{Questions: 1, 2, 3, 4}
    \item \textbf{C2: Undamped Resonance \& The $t$-Multiplier Trick} (Recognizing $\gamma = \omega$ when $b=0$ and multiplying the particular solution guess by $t$). \textit{Questions: 5, 6, 7, 8}
    \item \textbf{C3: The Beats Phenomenon} (Solving undamped systems where $\gamma \approx \omega$ and using sum-to-product trigonometric identities). \textit{Questions: 9, 10, 11}
    \item \textbf{C4: Steady-State Extraction} (Bypassing Undetermined Coefficients to directly find amplitude $A$ and phase angle $\theta$, including strict quadrant checking). \textit{Questions: 12, 13, 14, 15, 16}
    \item \textbf{C5: Full Initial Value Problems (Metric)} (Finding $y_h(t) + y_p(t)$ and applying initial conditions to the \textit{entire} system). \textit{Questions: 17, 18, 19, 20}
    \item \textbf{C6: Unit Conversions \& US Customary Systems} (Converting weight in lbs to mass in slugs using $m = W/32$, and finding $k$ via Hooke's Law). \textit{Questions: 21, 22, 23}
    \item \textbf{C7: The Complex Exponential Method} (Using Euler's formula $e^{i\gamma t}$ to bypass 2x2 systems of equations for steady-state solutions). \textit{Questions: 24, 25, 26}
    \item \textbf{C8: Advanced Diagnostic Testing ("Find the Flaw")} (Identifying subtle conceptual or algebraic errors in provided solutions). \textit{Questions: 27, 28, 29, 30}
\end{itemize}

\newpage

% ---------------------------------------------------------
% Problem Sets
% ---------------------------------------------------------
\section{Problem Sets}

\subsection*{Part 1: Frequency Response and Damped Resonance (C1)}
\noindent \textbf{Directions:} For the following systems governed by $my'' + by' + ky = F_0 \cos(\gamma t)$, determine if a resonance peak exists. If it does, find the resonance frequency $\gamma_r$ and the maximum frequency gain $M(\gamma_r)$.

\begin{enumerate}[label=\textbf{\arabic*.}]
    \item $m = 2$, $k = 8$, $b = 4$
    \item $m = 1$, $k = 9$, $b = 2$
    \item $m = 3$, $k = 27$, $b = 6$
    \item $m = 1$, $k = 4$, $b = 4$
\end{enumerate}

\subsection*{Part 2: Undamped Resonance (C2)}
\noindent \textbf{Directions:} Determine the particular solution $y_p(t)$ for the following undamped systems at resonance.

\begin{enumerate}[label=\textbf{\arabic*.}, resume]
    \item $y'' + 16y = 5 \cos(4t)$
    \item $y'' + 25y = 10 \cos(5t)$
    \item $2y'' + 18y = 3 \cos(3t)$
    \item $y'' + 36y = 12 \sin(6t)$
\end{enumerate}

\subsection*{Part 3: The Beats Phenomenon (C3)}
\noindent \textbf{Directions:} For the following undamped initial value problems, find the equation of motion $y(t)$. Then, use sum-to-product trigonometric identities to express the solution in a form that clearly shows the "beats" phenomenon.

\begin{enumerate}[label=\textbf{\arabic*.}, resume]
    \item $y'' + 100y = \cos(8t); \quad y(0) = 0, \quad y'(0) = 0$
    \item $y'' + 64y = 2\cos(6t); \quad y(0) = 0, \quad y'(0) = 0$
    \item $y'' + 144y = 4\cos(10t); \quad y(0) = 0, \quad y'(0) = 0$
\end{enumerate}

\subsection*{Part 4: Steady-State Extraction (C4)}
\noindent \textbf{Directions:} Find the steady-state solution $y_p(t) = A \sin(\gamma t + \theta)$ for the following systems. Be sure to correctly identify the quadrant for the phase angle $\theta$.

\begin{enumerate}[label=\textbf{\arabic*.}, resume]
    \item $y'' + 2y' + 5y = 10 \cos(t)$
    \item $y'' + 4y' + 3y = 5 \cos(2t)$
    \item $2y'' + 3y' + y = 4 \cos(3t)$
    \item $y'' + y' + 2y = 6 \cos(2t)$
    \item $0.5y'' + 2y' + 8y = 10 \cos(4t)$
\end{enumerate}

\subsection*{Part 5: Full Initial Value Problems (C5)}
\noindent \textbf{Directions:} Determine the full equation of motion $y(t)$ for the following systems.

\begin{enumerate}[label=\textbf{\arabic*.}, resume]
    \item $y'' + 2y' + 2y = 5 \cos(t); \quad y(0) = 1, \quad y'(0) = 0$
    \item $y'' + 4y' + 5y = 10 \cos(t); \quad y(0) = 0, \quad y'(0) = 0$
    \item $y'' + 3y' + 2y = 10 \sin(t); \quad y(0) = 0, \quad y'(0) = 0$
    \item $y'' + 6y' + 9y = 25 \cos(2t); \quad y(0) = 1, \quad y'(0) = -1$
\end{enumerate}

\subsection*{Part 6: US Customary Units \& Word Problems (C6)}
\noindent \textbf{Directions:} Set up the differential equation and find the steady-state solution $y_p(t)$. Pay strict attention to units.

\begin{enumerate}[label=\textbf{\arabic*.}, resume]
    \item A 16-lb weight is attached to a spring, stretching it $8/3$ ft to equilibrium. The damping constant is 2 lb-sec/ft. An external force $F(t) = 10 \cos(4t)$ is applied. 
    \item A 32-lb weight is attached to a spring with a spring constant of 17 lb/ft. The damping constant is 2 lb-sec/ft. An external force $F(t) = 4 \cos(4t)$ is applied.
    \item A 64-lb weight stretches a spring 2 ft to equilibrium. The damping constant is 4 lb-sec/ft. An external force $F(t) = 5 \cos(2t)$ is applied.
\end{enumerate}

\subsection*{Part 7: The Complex Exponential Method (C7)}
\noindent \textbf{Directions:} Use the Complex Exponential Method (Euler's Formula) to find the steady-state solution $y_p(t)$ for the following systems. Do not use the standard Method of Undetermined Coefficients.

\begin{enumerate}[label=\textbf{\arabic*.}, resume]
    \item $y'' + 2y' + 10y = 13 \cos(3t)$
    \item $y'' + 3y' + 2y = 10 \cos(2t)$
    \item $y'' + y' + 5y = 17 \cos(2t)$
\end{enumerate}

\subsection*{Part 8: Advanced Diagnostic Testing: "Find the Flaw" (C8)}
\noindent \textbf{Directions:} Each of the following solutions contains a single, critical flaw. Identify the error, explain why it is wrong, and provide the correct final answer.

\begin{enumerate}[label=\textbf{\arabic*.}, resume]
    \item \textbf{Problem:} Find the general solution to $y'' + 25y = 2 \cos(5t)$. \\
    \textbf{Flawed Solution:} 1. Homogeneous: $r = \pm 5i$. 2. Particular: Guess $y_p = A \cos(5t) + B \sin(5t)$. 3. Differentiate and substitute: $-25A \cos(5t) - 25B \sin(5t) + 25A \cos(5t) + 25B \sin(5t) = 2 \cos(5t)$. 4. Simplify: $0 = 2 \cos(5t)$. Conclusion: No solution exists.
    
    \item \textbf{Problem:} A 32-lb weight is attached to a spring with $k = 4$ lb/ft. Damping is $b = 2$. Find the resonance frequency $\gamma_r$. \\
    \textbf{Flawed Solution:} 1. Parameters: $m = 32, k = 4, b = 2$. 2. Formula: $\gamma_r = \sqrt{\frac{k}{m} - \frac{b^2}{2m^2}}$. 3. Substitute: $\gamma_r = \sqrt{\frac{4}{32} - \frac{4}{2(1024)}} = \sqrt{\frac{1}{8} - \frac{4}{2048}} \approx 0.35$ rad/sec.
    
    \item \textbf{Problem:} Find the phase angle $\theta$ for the steady-state solution of $y'' + 2y' + y = \cos(2t)$. \\
    \textbf{Flawed Solution:} 1. Parameters: $m=1, b=2, k=1, \gamma=2$. 2. Formula: $\tan \theta = \frac{k-m\gamma^2}{b\gamma}$. 3. Substitute: $\tan \theta = \frac{1 - 1(4)}{2(2)} = -\frac{3}{4}$. 4. Calculate: $\theta = \arctan(-0.75) \approx -0.64$ radians. 5. Because the tangent is negative, we must add $\pi$ to find the true angle. 6. Final Answer: $\theta = -0.64 + \pi = 2.50$ radians (Quadrant II).
    
    \item \textbf{Problem:} Find the amplitude of the steady-state solution for $y'' + 3y' + 4y = 10 \cos(t)$. \\
    \textbf{Flawed Solution:} 1. Parameters: $m=1, b=3, k=4, F_0=10, \gamma=1$. 2. Formula: $A = \frac{F_0}{\sqrt{(k-m\gamma^2)^2 + (b\gamma)^2}}$. 3. Substitute: $A = \frac{10}{\sqrt{(4-1(1))^2 + (3(1))^2}}$. 4. Simplify: $A = \frac{10}{\sqrt{3^2 + 3^2}} = \frac{10}{3 + 3} = \frac{10}{6} = 1.67$.
\end{enumerate}

\newpage

% ---------------------------------------------------------
% Detailed Solutions
% ---------------------------------------------------------
\section{Detailed Solutions}

\subsection*{Part 1: Frequency Response and Damped Resonance}

\textbf{1.} $m = 2, k = 8, b = 4$. \\
Check for resonance: $b^2 = 16$. $2mk = 2(2)(8) = 32$. Since $16 < 32$, a resonance peak exists. \\
$\gamma_r = \sqrt{\frac{k}{m} - \frac{b^2}{2m^2}} = \sqrt{\frac{8}{2} - \frac{16}{2(4)}} = \sqrt{4 - 2} = \sqrt{2}$ rad/sec. \\
Maximum gain: $M(\sqrt{2}) = \frac{1}{\sqrt{(8 - 2(2))^2 + 4^2(2)}} = \frac{1}{\sqrt{16 + 32}} = \frac{1}{\sqrt{48}} = \frac{1}{4\sqrt{3}}$.

\vspace{0.5cm}
\textbf{2.} $m = 1, k = 9, b = 2$. \\
Check for resonance: $b^2 = 4$. $2mk = 2(1)(9) = 18$. Since $4 < 18$, a resonance peak exists. \\
$\gamma_r = \sqrt{\frac{9}{1} - \frac{4}{2(1)}} = \sqrt{9 - 2} = \sqrt{7}$ rad/sec. \\
Maximum gain: $M(\sqrt{7}) = \frac{1}{\sqrt{(9 - 1(7))^2 + 2^2(7)}} = \frac{1}{\sqrt{4 + 28}} = \frac{1}{\sqrt{32}} = \frac{1}{4\sqrt{2}}$.

\vspace{0.5cm}
\textbf{3.} $m = 3, k = 27, b = 6$. \\
Check for resonance: $b^2 = 36$. $2mk = 2(3)(27) = 162$. Since $36 < 162$, a resonance peak exists. \\
$\gamma_r = \sqrt{\frac{27}{3} - \frac{36}{2(9)}} = \sqrt{9 - 2} = \sqrt{7}$ rad/sec. \\
Maximum gain: $M(\sqrt{7}) = \frac{1}{\sqrt{(27 - 3(7))^2 + 6^2(7)}} = \frac{1}{\sqrt{36 + 252}} = \frac{1}{\sqrt{288}} = \frac{1}{12\sqrt{2}}$.

\vspace{0.5cm}
\textbf{4.} $m = 1, k = 4, b = 4$. \\
Check for resonance: $b^2 = 16$. $2mk = 2(1)(4) = 8$. Since $16 \ge 8$, the system is too damped to resonate. There is no resonance peak (maximum amplitude occurs at $\gamma = 0$).

\subsection*{Part 2: Undamped Resonance}

\textbf{5.} $y'' + 16y = 5 \cos(4t)$. \\
Natural frequency $\omega = \sqrt{16} = 4$. Forcing frequency $\gamma = 4$. Since they match, guess $y_p = t(A \cos 4t + B \sin 4t)$. \\
$y_p'' = -16A t \cos 4t - 16B t \sin 4t - 8A \sin 4t + 8B \cos 4t$. \\
Substitute into ODE: $(-8A \sin 4t + 8B \cos 4t) = 5 \cos 4t$. \\
$A = 0, 8B = 5 \implies B = \frac{5}{8}$. \\
\textbf{Answer:} $y_p(t) = \frac{5}{8} t \sin(4t)$.

\vspace{0.5cm}
\textbf{6.} $y'' + 25y = 10 \cos(5t)$. \\
$\omega = 5, \gamma = 5$. Guess $y_p = t(A \cos 5t + B \sin 5t)$. \\
$y_p'' + 25y_p = -10A \sin 5t + 10B \cos 5t = 10 \cos 5t$. \\
$A = 0, 10B = 10 \implies B = 1$. \\
\textbf{Answer:} $y_p(t) = t \sin(5t)$.

\vspace{0.5cm}
\textbf{7.} $2y'' + 18y = 3 \cos(3t) \implies y'' + 9y = 1.5 \cos(3t)$. \\
$\omega = 3, \gamma = 3$. Guess $y_p = t(A \cos 3t + B \sin 3t)$. \\
$y_p'' + 9y_p = -6A \sin 3t + 6B \cos 3t = 1.5 \cos 3t$. \\
$A = 0, 6B = 1.5 \implies B = \frac{1.5}{6} = \frac{1}{4}$. \\
\textbf{Answer:} $y_p(t) = \frac{1}{4} t \sin(3t)$.

\vspace{0.5cm}
\textbf{8.} $y'' + 36y = 12 \sin(6t)$. \\
$\omega = 6, \gamma = 6$. Guess $y_p = t(A \cos 6t + B \sin 6t)$. \\
$y_p'' + 36y_p = -12A \sin 6t + 12B \cos 6t = 12 \sin 6t$. \\
$-12A = 12 \implies A = -1$. $12B = 0 \implies B = 0$. \\
\textbf{Answer:} $y_p(t) = -t \cos(6t)$.

\subsection*{Part 3: The Beats Phenomenon}

\textbf{9.} $y'' + 100y = \cos(8t); \quad y(0) = 0, y'(0) = 0$. \\
$y_h = c_1 \cos(10t) + c_2 \sin(10t)$. \\
$y_p = A \cos(8t) \implies A(-64) + 100A = 1 \implies 36A = 1 \implies A = \frac{1}{36}$. \\
$y(t) = c_1 \cos(10t) + c_2 \sin(10t) + \frac{1}{36} \cos(8t)$. \\
$y(0) = c_1 + \frac{1}{36} = 0 \implies c_1 = -\frac{1}{36}$. \\
$y'(0) = 10c_2 = 0 \implies c_2 = 0$. \\
$y(t) = \frac{1}{36}(\cos 8t - \cos 10t)$. \\
Sum-to-product: $\cos A - \cos B = -2 \sin(\frac{A+B}{2}) \sin(\frac{A-B}{2})$. \\
$y(t) = \frac{1}{36} \left[ -2 \sin(9t) \sin(-t) \right] = \frac{1}{18} \sin(t) \sin(9t)$.

\vspace{0.5cm}
\textbf{10.} $y'' + 64y = 2\cos(6t); \quad y(0) = 0, y'(0) = 0$. \\
$y_h = c_1 \cos(8t) + c_2 \sin(8t)$. \\
$y_p = A \cos(6t) \implies A(-36) + 64A = 2 \implies 28A = 2 \implies A = \frac{1}{14}$. \\
$y(0) = c_1 + \frac{1}{14} = 0 \implies c_1 = -\frac{1}{14}, c_2 = 0$. \\
$y(t) = \frac{1}{14}(\cos 6t - \cos 8t) = \frac{1}{14} \left[ -2 \sin(7t) \sin(-t) \right] = \frac{1}{7} \sin(t) \sin(7t)$.

\vspace{0.5cm}
\textbf{11.} $y'' + 144y = 4\cos(10t); \quad y(0) = 0, y'(0) = 0$. \\
$y_h = c_1 \cos(12t) + c_2 \sin(12t)$. \\
$y_p = A \cos(10t) \implies A(-100) + 144A = 4 \implies 44A = 4 \implies A = \frac{1}{11}$. \\
$y(0) = c_1 + \frac{1}{11} = 0 \implies c_1 = -\frac{1}{11}, c_2 = 0$. \\
$y(t) = \frac{1}{11}(\cos 10t - \cos 12t) = \frac{1}{11} \left[ -2 \sin(11t) \sin(-t) \right] = \frac{2}{11} \sin(t) \sin(11t)$.

\subsection*{Part 4: Steady-State Extraction}

\textbf{12.} $y'' + 2y' + 5y = 10 \cos(t)$. \\
$m=1, b=2, k=5, F_0=10, \gamma=1$. \\
$k - m\gamma^2 = 5 - 1(1) = 4$. $b\gamma = 2(1) = 2$. \\
Amplitude $A = \frac{10}{\sqrt{4^2 + 2^2}} = \frac{10}{\sqrt{20}} = \frac{10}{2\sqrt{5}} = \sqrt{5}$. \\
Phase $\tan \theta = \frac{4}{2} = 2$. Numerator (+), Denominator (+) $\implies$ Quadrant I. \\
\textbf{Answer:} $y_p(t) = \sqrt{5} \sin(t + \arctan(2))$.

\vspace{0.5cm}
\textbf{13.} $y'' + 4y' + 3y = 5 \cos(2t)$. \\
$m=1, b=4, k=3, F_0=5, \gamma=2$. \\
$k - m\gamma^2 = 3 - 1(4) = -1$. $b\gamma = 4(2) = 8$. \\
Amplitude $A = \frac{5}{\sqrt{(-1)^2 + 8^2}} = \frac{5}{\sqrt{65}}$. \\
Phase $\tan \theta = -\frac{1}{8}$. Numerator (-), Denominator (+) $\implies$ Quadrant IV. \\
\textbf{Answer:} $y_p(t) = \frac{5}{\sqrt{65}} \sin(2t + \arctan(-1/8))$.

\vspace{0.5cm}
\textbf{14.} $2y'' + 3y' + y = 4 \cos(3t)$. \\
$m=2, b=3, k=1, F_0=4, \gamma=3$. \\
$k - m\gamma^2 = 1 - 2(9) = -17$. $b\gamma = 3(3) = 9$. \\
Amplitude $A = \frac{4}{\sqrt{(-17)^2 + 9^2}} = \frac{4}{\sqrt{289 + 81}} = \frac{4}{\sqrt{370}}$. \\
Phase $\tan \theta = -\frac{17}{9}$. Numerator (-), Denominator (+) $\implies$ Quadrant IV. \\
\textbf{Answer:} $y_p(t) = \frac{4}{\sqrt{370}} \sin(3t + \arctan(-17/9))$.

\vspace{0.5cm}
\textbf{15.} $y'' + y' + 2y = 6 \cos(2t)$. \\
$m=1, b=1, k=2, F_0=6, \gamma=2$. \\
$k - m\gamma^2 = 2 - 1(4) = -2$. $b\gamma = 1(2) = 2$. \\
Amplitude $A = \frac{6}{\sqrt{(-2)^2 + 2^2}} = \frac{6}{\sqrt{8}} = \frac{6}{2\sqrt{2}} = \frac{3\sqrt{2}}{2}$. \\
Phase $\tan \theta = -\frac{2}{2} = -1$. Numerator (-), Denominator (+) $\implies$ Quadrant IV. $\theta = -\frac{\pi}{4}$. \\
\textbf{Answer:} $y_p(t) = \frac{3\sqrt{2}}{2} \sin(2t - \frac{\pi}{4})$.

\vspace{0.5cm}
\textbf{16.} $0.5y'' + 2y' + 8y = 10 \cos(4t)$. \\
$m=0.5, b=2, k=8, F_0=10, \gamma=4$. \\
$k - m\gamma^2 = 8 - 0.5(16) = 0$. $b\gamma = 2(4) = 8$. \\
Amplitude $A = \frac{10}{\sqrt{0^2 + 8^2}} = \frac{10}{8} = 1.25$. \\
Phase $\tan \theta = \frac{0}{8} = 0 \implies \theta = 0$. \\
\textbf{Answer:} $y_p(t) = 1.25 \sin(4t)$.

\subsection*{Part 5: Full Initial Value Problems}

\textbf{17.} $y'' + 2y' + 2y = 5 \cos(t); \quad y(0) = 1, y'(0) = 0$. \\
Homogeneous: $r^2 + 2r + 2 = 0 \implies r = -1 \pm i$. $y_h = e^{-t}(c_1 \cos t + c_2 \sin t)$. \\
Particular: Guess $y_p = A \cos t + B \sin t$. \\
$y_p' = -A \sin t + B \cos t$, $y_p'' = -A \cos t - B \sin t$. \\
Substitute: $(-A + 2B + 2A)\cos t + (-B - 2A + 2B)\sin t = 5 \cos t$. \\
$A + 2B = 5$ and $-2A + B = 0 \implies B = 2A \implies 5A = 5 \implies A=1, B=2$. \\
$y_p = \cos t + 2 \sin t$. \\
General Solution: $y(t) = e^{-t}(c_1 \cos t + c_2 \sin t) + \cos t + 2 \sin t$. \\
$y(0) = c_1 + 1 = 1 \implies c_1 = 0$. \\
$y'(t) = e^{-t}(-c_2 \sin t) + e^{-t}(c_2 \cos t) - \sin t + 2 \cos t$. \\
$y'(0) = c_2 + 2 = 0 \implies c_2 = -2$. \\
\textbf{Answer:} $y(t) = -2e^{-t}\sin t + \cos t + 2 \sin t$.

\vspace{0.5cm}
\textbf{18.} $y'' + 4y' + 5y = 10 \cos(t); \quad y(0) = 0, y'(0) = 0$. \\
Homogeneous: $r^2 + 4r + 5 = 0 \implies r = -2 \pm i$. $y_h = e^{-2t}(c_1 \cos t + c_2 \sin t)$. \\
Particular: Guess $y_p = A \cos t + B \sin t$. \\
Substitute: $(-A + 4B + 5A)\cos t + (-B - 4A + 5B)\sin t = 10 \cos t$. \\
$4A + 4B = 10$ and $-4A + 4B = 0 \implies A = B \implies 8A = 10 \implies A = 1.25, B = 1.25$. \\
$y_p = 1.25 \cos t + 1.25 \sin t$. \\
General Solution: $y(t) = e^{-2t}(c_1 \cos t + c_2 \sin t) + 1.25 \cos t + 1.25 \sin t$. \\
$y(0) = c_1 + 1.25 = 0 \implies c_1 = -1.25$. \\
$y'(t) = -2e^{-2t}(c_1 \cos t + c_2 \sin t) + e^{-2t}(-c_1 \sin t + c_2 \cos t) - 1.25 \sin t + 1.25 \cos t$. \\
$y'(0) = -2c_1 + c_2 + 1.25 = 0 \implies 2.5 + c_2 + 1.25 = 0 \implies c_2 = -3.75$. \\
\textbf{Answer:} $y(t) = e^{-2t}(-1.25 \cos t - 3.75 \sin t) + 1.25 \cos t + 1.25 \sin t$.

\vspace{0.5cm}
\textbf{19.} $y'' + 3y' + 2y = 10 \sin(t); \quad y(0) = 0, y'(0) = 0$. \\
Homogeneous: $r^2 + 3r + 2 = 0 \implies r = -1, -2$. $y_h = c_1 e^{-t} + c_2 e^{-2t}$. \\
Particular: Guess $y_p = A \cos t + B \sin t$. \\
Substitute: $(-A + 3B + 2A)\cos t + (-B - 3A + 2B)\sin t = 10 \sin t$. \\
$A + 3B = 0 \implies A = -3B$. $-3A + B = 10 \implies 9B + B = 10 \implies B = 1, A = -3$. \\
$y_p = -3 \cos t + \sin t$. \\
General Solution: $y(t) = c_1 e^{-t} + c_2 e^{-2t} - 3 \cos t + \sin t$. \\
$y(0) = c_1 + c_2 - 3 = 0 \implies c_1 + c_2 = 3$. \\
$y'(0) = -c_1 - 2c_2 + 1 = 0 \implies c_1 + 2c_2 = 1$. \\
Subtracting equations: $c_2 = -2 \implies c_1 = 5$. \\
\textbf{Answer:} $y(t) = 5e^{-t} - 2e^{-2t} - 3 \cos t + \sin t$.

\vspace{0.5cm}
\textbf{20.} $y'' + 6y' + 9y = 25 \cos(2t); \quad y(0) = 1, y'(0) = -1$. \\
Homogeneous: $r^2 + 6r + 9 = 0 \implies r = -3$ (double root). $y_h = c_1 e^{-3t} + c_2 t e^{-3t}$. \\
Particular: Guess $y_p = A \cos 2t + B \sin 2t$. \\
Substitute: $(-4A + 12B + 9A)\cos 2t + (-4B - 12A + 9B)\sin 2t = 25 \cos 2t$. \\
$5A + 12B = 25$ and $-12A + 5B = 0 \implies B = \frac{12}{5}A$. \\
$5A + \frac{144}{5}A = 25 \implies \frac{169}{5}A = 25 \implies A = \frac{125}{169}, B = \frac{300}{169}$. \\
General Solution: $y(t) = c_1 e^{-3t} + c_2 t e^{-3t} + \frac{125}{169} \cos 2t + \frac{300}{169} \sin 2t$. \\
$y(0) = c_1 + \frac{125}{169} = 1 \implies c_1 = \frac{44}{169}$. \\
$y'(0) = -3c_1 + c_2 + \frac{600}{169} = -1 \implies -\frac{132}{169} + c_2 + \frac{600}{169} = -\frac{169}{169} \implies c_2 = -\frac{637}{169} = -\frac{49}{13}$. \\
\textbf{Answer:} $y(t) = \frac{44}{169} e^{-3t} - \frac{49}{13} t e^{-3t} + \frac{125}{169} \cos 2t + \frac{300}{169} \sin 2t$.

\subsection*{Part 6: US Customary Units \& Word Problems}

\textbf{21.} Mass $m = \frac{W}{g} = \frac{16}{32} = 0.5$ slugs. \\
Spring constant $k = \frac{mg}{L} = \frac{16}{8/3} = 6$ lb/ft. \\
$b = 2, F_0 = 10, \gamma = 4$. \\
$k - m\gamma^2 = 6 - 0.5(16) = -2$. $b\gamma = 2(4) = 8$. \\
Amplitude $A = \frac{10}{\sqrt{(-2)^2 + 8^2}} = \frac{10}{\sqrt{68}} = \frac{5}{\sqrt{17}}$. \\
Phase $\tan \theta = -\frac{2}{8} = -0.25$. Quadrant IV. \\
\textbf{Answer:} $y_p(t) = \frac{5}{\sqrt{17}} \sin(4t + \arctan(-0.25))$.

\vspace{0.5cm}
\textbf{22.} Mass $m = \frac{32}{32} = 1$ slug. \\
$k = 17, b = 2, F_0 = 4, \gamma = 4$. \\
$k - m\gamma^2 = 17 - 1(16) = 1$. $b\gamma = 2(4) = 8$. \\
Amplitude $A = \frac{4}{\sqrt{1^2 + 8^2}} = \frac{4}{\sqrt{65}}$. \\
Phase $\tan \theta = \frac{1}{8}$. Quadrant I. \\
\textbf{Answer:} $y_p(t) = \frac{4}{\sqrt{65}} \sin(4t + \arctan(1/8))$.

\vspace{0.5cm}
\textbf{23.} Mass $m = \frac{64}{32} = 2$ slugs. \\
Spring constant $k = \frac{64}{2} = 32$ lb/ft. \\
$b = 4, F_0 = 5, \gamma = 2$. \\
$k - m\gamma^2 = 32 - 2(4) = 24$. $b\gamma = 4(2) = 8$. \\
Amplitude $A = \frac{5}{\sqrt{24^2 + 8^2}} = \frac{5}{\sqrt{576 + 64}} = \frac{5}{\sqrt{640}} = \frac{5}{8\sqrt{10}}$. \\
Phase $\tan \theta = \frac{24}{8} = 3$. Quadrant I. \\
\textbf{Answer:} $y_p(t) = \frac{5}{8\sqrt{10}} \sin(2t + \arctan(3))$.

\subsection*{Part 7: The Complex Exponential Method}

\textbf{24.} $y'' + 2y' + 10y = 13 \cos(3t)$. \\
Complex ODE: $z'' + 2z' + 10z = 13 e^{3it}$. Guess $z_p = C e^{3it}$. \\
Substitute: $C(-9 + 6i + 10) = 13 \implies C(1 + 6i) = 13$. \\
$C = \frac{13}{1 + 6i} = \frac{13(1 - 6i)}{1 + 36} = \frac{13(1 - 6i)}{37} = \frac{13}{37} - \frac{78}{37}i$. \\
$z_p = (\frac{13}{37} - \frac{78}{37}i)(\cos 3t + i \sin 3t)$. \\
Take the real part: $y_p(t) = \text{Re}(z_p) = \frac{13}{37} \cos 3t + \frac{78}{37} \sin 3t$.

\vspace{0.5cm}
\textbf{25.} $y'' + 3y' + 2y = 10 \cos(2t)$. \\
Complex ODE: $z'' + 3z' + 2z = 10 e^{2it}$. Guess $z_p = C e^{2it}$. \\
Substitute: $C(-4 + 6i + 2) = 10 \implies C(-2 + 6i) = 10$. \\
$C = \frac{10}{-2 + 6i} = \frac{10(-2 - 6i)}{4 + 36} = \frac{-20 - 60i}{40} = -0.5 - 1.5i$. \\
$z_p = (-0.5 - 1.5i)(\cos 2t + i \sin 2t)$. \\
Take the real part: $y_p(t) = \text{Re}(z_p) = -0.5 \cos 2t + 1.5 \sin 2t$.

\vspace{0.5cm}
\textbf{26.} $y'' + y' + 5y = 17 \cos(2t)$. \\
Complex ODE: $z'' + z' + 5z = 17 e^{2it}$. Guess $z_p = C e^{2it}$. \\
Substitute: $C(-4 + 2i + 5) = 17 \implies C(1 + 2i) = 17$. \\
$C = \frac{17}{1 + 2i} = \frac{17(1 - 2i)}{1 + 4} = \frac{17(1 - 2i)}{5} = 3.4 - 6.8i$. \\
$z_p = (3.4 - 6.8i)(\cos 2t + i \sin 2t)$. \\
Take the real part: $y_p(t) = \text{Re}(z_p) = 3.4 \cos 2t + 6.8 \sin 2t$.

\subsection*{Part 8: Advanced Diagnostic Testing: "Find the Flaw"}

\textbf{27.} \textbf{Flaw:} The forcing frequency ($\gamma = 5$) exactly matches the natural frequency ($\omega = \sqrt{25} = 5$) of the undamped system. This is resonance, so the standard guess fails and must be multiplied by $t$. \\
\textbf{Correct Answer:} $y_p(t) = \frac{1}{5} t \sin(5t)$.

\vspace{0.5cm}
\textbf{28.} \textbf{Flaw:} The weight (32 lbs) was incorrectly used as the mass $m$. In the US Customary System, mass must be calculated as $m = W/g = 32/32 = 1$ slug. \\
\textbf{Correct Answer:} $\gamma_r = \sqrt{\frac{4}{1} - \frac{4}{2(1)}} = \sqrt{4 - 2} = \sqrt{2}$ rad/sec.

\vspace{0.5cm}
\textbf{29.} \textbf{Flaw:} The denominator $b\gamma = 4$ is positive, and the numerator $k-m\gamma^2 = -3$ is negative. This places the angle strictly in Quadrant IV. Adding $\pi$ incorrectly shifts the angle to Quadrant II. \\
\textbf{Correct Answer:} $\theta = \arctan(-0.75) \approx -0.64$ radians.

\vspace{0.5cm}
\textbf{30.} \textbf{Flaw:} The algebraic simplification $\sqrt{3^2 + 3^2} = 3 + 3$ is a severe violation of order of operations. The square root of a sum of squares does not equal the sum of the square roots. \\
\textbf{Correct Answer:} $A = \frac{10}{\sqrt{9 + 9}} = \frac{10}{\sqrt{18}} = \frac{10}{3\sqrt{2}} = \frac{5\sqrt{2}}{3}$.

\newpage

% ---------------------------------------------------------
% Appendix: Concept Checklist & Question Mapping (Repeated)
% ---------------------------------------------------------
\section*{Appendix: Concept Checklist \& Question Mapping}
\addcontentsline{toc}{section}{Appendix: Concept Checklist \& Question Mapping}

\begin{itemize}
    \item \textbf{C1: Frequency Response Curve \& Damped Resonance} (Checking $b^2 < 2mk$, calculating $\gamma_r$ and maximum gain $M(\gamma_r)$). \textit{Questions: 1, 2, 3, 4}
    \item \textbf{C2: Undamped Resonance \& The $t$-Multiplier Trick} (Recognizing $\gamma = \omega$ when $b=0$ and multiplying the particular solution guess by $t$). \textit{Questions: 5, 6, 7, 8}
    \item \textbf{C3: The Beats Phenomenon} (Solving undamped systems where $\gamma \approx \omega$ and using sum-to-product trigonometric identities). \textit{Questions: 9, 10, 11}
    \item \textbf{C4: Steady-State Extraction} (Bypassing Undetermined Coefficients to directly find amplitude $A$ and phase angle $\theta$, including strict quadrant checking). \textit{Questions: 12, 13, 14, 15, 16}
    \item \textbf{C5: Full Initial Value Problems (Metric)} (Finding $y_h(t) + y_p(t)$ and applying initial conditions to the \textit{entire} system). \textit{Questions: 17, 18, 19, 20}
    \item \textbf{C6: Unit Conversions \& US Customary Systems} (Converting weight in lbs to mass in slugs using $m = W/32$, and finding $k$ via Hooke's Law). \textit{Questions: 21, 22, 23}
    \item \textbf{C7: The Complex Exponential Method} (Using Euler's formula $e^{i\gamma t}$ to bypass 2x2 systems of equations for steady-state solutions). \textit{Questions: 24, 25, 26}
    \item \textbf{C8: Advanced Diagnostic Testing ("Find the Flaw")} (Identifying subtle conceptual or algebraic errors in provided solutions). \textit{Questions: 27, 28, 29, 30}
\end{itemize}

\end{document}